\documentclass[a4paper,11pt]{article}% {{{
% -------------------------------------------------
% {{{ Set up the page margins.
% -------------------------------------------------
\usepackage[margin=1in]{geometry}
% \usepackage[headheig\label{key}ht = .5in, headsep = \baselineskip, top = 1in, left = 1in, right = 1in, textwidth = 7.52 in]{geometry}
% \usepackage[text={6.5in,9.5in},centering]{geometry}
% \usepackage[margin=1in]{geometry}
% \setlength{\topmargin}{-0.25in}
% }}}
% -------------------------------------------------
% {{{ Load Packages here
% -------------------------------------------------
\usepackage[utf8]{inputenc}
\usepackage[T1]{fontenc}
\usepackage{amssymb,
  nccmath,
  geometry,
    latexsym,
    amsmath,
    graphics,
    fullpage,
    epsfig,
    amsthm,
    relsize,
    tikz,
    amsfonts,
    mathrsfs,
    makeidx,
    comment,
    latexsym,
    dcolumn,
    tocloft, % spacing for tableofcontents
    calc}
\usepackage[colorlinks,
     citecolor=blue,
     urlcolor=blue
      ]{hyperref}
\usepackage[shortlabels]{enumitem}
\excludecomment{codes}
% }}}
% -------------------------------------------------
% {{{ check references
% -------------------------------------------------
\usepackage{refcheck}
% \usepackage[notref,notcite]{showkeys}
% \usepackage[notcite]{showkeys}
% \usepackage{showkeys}
% }}}
% -------------------------------------------------
% {{{ Biber related
% -------------------------------------------------
\usepackage[strict=true,
            style=english]{csquotes}
% \usepackage{mathscinet}
\DeclareTextCommand{\polhk}{OT4}{\k}
\DeclareTextCommand{\polhk}{T1}{\k}

% \setlength{\parskip}{0.5cm plus4mm minus3mm}
% \usepackage[style=alphabetic, %authoryear
%             bibstyle=authoryear,
%             citestyle=authoryear,
%             natbib=true,
%             hyperref=true,
%             backref=true,
%             abbreviate=true]{biblatex}

% \addbibresource{./draft_rand_int_biber.bib}
% }}}
% -------------------------------------------------
% {{{ Environments and some setups
% -------------------------------------------------
\theoremstyle{plain}
\newtheorem{theorem}{Theorem}[section]
\newtheorem{corollary}[theorem]{Corollary}
\newtheorem{lemma}[theorem]{Lemma}
\newtheorem{proposition}[theorem]{Proposition}

\theoremstyle{definition}
\newtheorem{definition}[theorem]{Definition}
\newtheorem{hypothesis}[theorem]{Hypothesis}
\newtheorem{assumptions}[theorem]{Assumption}

\newtheorem{remark}[theorem]{Remark}
\newtheorem{conjecture}[theorem]{Conjecture}
\newtheorem{example}[theorem]{Example}
\newtheorem{notation}[theorem]{Notation}

\numberwithin{equation}{section}
\allowdisplaybreaks
% }}}
% -------------------------------------------------
% {{{ Biber related
% -------------------------------------------------
\usepackage[strict = true,
            style = english]{csquotes}
% \usepackage{mathscinet}
\DeclareTextCommand{\polhk}{OT4}{\k}
\DeclareTextCommand{\polhk}{T1}{\k}
\usepackage[backend = biber,
						% style = alphabetic,
						style = numeric,
						natbib = true,
            sorting = nyt,
            maxbibnames=99,
            abbreviate = true,
           ]{biblatex}
\AtEveryBibitem{% Clean up the bibtex rather than editing it
 \clearfield{doi}
 \clearfield{url}
 % \clearfield{issn}
 % \clearlist{location}
 \clearfield{month}
 \clearfield{series}

 \ifentrytype{book}{}{% Remove publisher and editor except for books
  \clearlist{publisher}
  \clearname{editor}
 }
}
\renewcommand*{\bibfont}{\fontsize{11}{11}}

\addbibresource{../All.bib}
% \addbibresource{./draft_rand_int_2_biber.bib}
% }}}
% -------------------------------------------------
% {{{ Define short hand symbols.
% -------------------------------------------------
\newcommand{\lip}{\textup{Lip}_b}
\newcommand{\liph}{\text{Lip}_{\widehat{\rho} }}
\newcommand{\R}{\mathbb{R}}
\newcommand{\E}{\mathbb{E}}
\newcommand{\e}{\mathrm{e}}
\newcommand{\Norm}[1]{\left\|  #1   \right\|}
\newcommand{\rhoNorm}[1]{\left|  #1   \right|}
\newcommand{\ud}{\ensuremath{ \mathrm{d} }}
\newcommand*{\one}{ {{\rm 1\mkern-1.5mu}\!{\rm I} }}
% }}}
% }}}
% {{{ Head: Title, Abstract, Keywords
\title{Comments on the reviewing report for\\
Invariant measures for the stochastic heat equation on $\R^d$ with no drift term}
\author{Le Chen\\Auburn University\\\texttt{le.chen@auburn.edu}\\
  \and
    Nicholas Eisenberg\\Auburn University\\\texttt{nze0019@auburn.edu}
  }
\date{\today}

\begin{document}
  \maketitle
 
  We thank the anonymous referee for careful reading of the paper, some helpful
  comments, and especially pointing out the reference by Gu and
  Li~\cite{gu.li:20:fluctuations}. In this revision, we have carefully taken
  into all comments from the report. The list of major changes are as follows:

	\begin{enumerate}
    \item pages 3 -- 4, in the both paragraphs after Eq.~(1.9), we have included
      the reference~\cite{gu.li:20:fluctuations} and made some comparison with
      our current work. As a short summary, the focus of our paper is on the
      conditions on the correlation function/measure and initial conditions in
      order to guarantee the existence of invariant measures. The focus of the
      paper~\cite{gu.li:20:fluctuations} is different, and convenient conditions
      were made in~\cite{gu.li:20:fluctuations} for the correlation function,
      which is assumed to be smooth function with compact support. The initial
      condition in~\cite{gu.li:20:fluctuations} is essentially constant initial
      data with some perturbation. In contrast, our conditions on both
      correlation function/measure and initial conditions are much weaker; see
      the main theorem -- Theorem 1.5 of the paper.
    \item page 5, Figure 1 has been simplified and updated. We postponed the
      conditions formulated using $\mathcal{H}_\alpha(t)$ back in Section 3.1
      when it is first used.
		\item page 5, Remark 1.2 has been added to comment on the including the
      \textit{strengthened Dalang's condition} in our main theorem.
		\item page 7, Remark 1.6 has been included to comment on our conditions on
      the noise. Both this change and the previous one are related to one
      comment from the referee about our additional assumption, namely, the
      \textit{strengthened Dalang's condition}.
		\item Section 3.2, Figure 2 and Remark 3.5 have been created to help the reader. 
      We have also updated Lemma 3.6 accordingly to reflect the changes in items
      2, 3, and 4 above.
    \item In the Acknowledgment section, we include acknowledgment to Yu Gu
      for the discussions on the reference~\cite{gu.li:20:fluctuations}.
	\end{enumerate}


% {{{All references: New ones with biber
\addcontentsline{toc}{section}{References}
\printbibliography[title = {References}]
%}}}
\end{document}
